\section{USB Power Delivery}\label{sec:usb-pd}

The USB-C input supplies all three converters. A CYPD3176 controller negotiates
the power contract. A pass switch connects the accepted supply to the board.
An analog switch controls the MCU USB data path. \Cref{tab:pd-designators}
defines the designators.

\begin{table}[htbp]
\centering
\caption{USB-PD front end: component designators}\label{tab:pd-designators}
\begin{tabular}{|l|l|}
\hline
\textbf{Designator} & \textbf{Function} \\
\hline
\des{J}{USB} & USB Type-C receptacle \\
\hline
\des{Q}{BUS} & Bus pass switch, two common-source devices driven by VBUS\_CTRL \\
\hline
\des{R}{BUSG} & \des{Q}{BUS} gate series resistor \\
\hline
\des{R}{BUSOFF} & \des{Q}{BUS} gate-to-source default-off pull-up \\
\hline
\des{R}{BUSBIAS} & \des{Q}{BUS} gate bias to the receptacle's VBUS \\
\hline
\des{D}{BUSG} & \des{Q}{BUS} gate-source clamp \\
\hline
\des{C}{BUSSLEW} & \des{Q}{BUS} output slew-rate capacitor \\
\hline
\des{C}{BUSCM} & \des{Q}{BUS} gate common-mode capacitor \\
\hline
\des{R}{PDMAX1}, \des{R}{PDMAX2} & VBUS\_MAX selection divider \\
\hline
\des{R}{PDI1}, \des{R}{PDI2} & ISNK\_FINE selection divider \\
\hline
\des{C}{PDD1} & VDDD bulk decoupling \\
\hline
\des{C}{PDD2}, \des{C}{PDD3} & VDDD high-frequency decoupling \\
\hline
\des{C}{PDC} & VCCD decoupling \\
\hline
\des{C}{PDIN} & VBUS\_IN\_DIS decoupling \\
\hline
\des{R}{PDSDA}, \des{R}{PDSCL} & HPI bus pull-ups \\
\hline
\des{R}{PDINT} & HPI\_INT pull-up \\
\hline
\des{R}{PDFLT} & FAULT mode-select pull-down \\
\hline
\des{R}{USBSW} & Data-switch disable pull-up \\
\hline
\des{R}{PDDP}, \des{R}{PDDM} & Controller D+/D$-$ branch series resistors \\
\hline
\des{D}{USBESD} & D+/D$-$ ESD array \\
\hline
\des{D}{BUSTVS} & Receptacle VBUS transient suppressor \\
\hline
\des{R}{BUSBLEED} & Bus-rail discharge path (\cref{sec:3v3-diodeor}) \\
\hline
\end{tabular}
\end{table}

\begin{figure}[htbp]
\centering
\includegraphics[width=\textwidth]{usbpd}
\caption{USB-PD front end: reference schematic}\label{fig:usbpd}
\end{figure}

\Cref{fig:usbpd} shows the connections. The labeled signal nets connect to the
MCU (\cref{sec:esp32}). $V_{\mathrm{BUS}}$ supplies the three converters.

\subsection{Driving component characteristics}\label{sec:cypd3176}

Source: \texttt{infineon-cypd3176-24lqxq-datasheet-en.pdf} (Infineon,
002-32090 Rev.~*E, 2024-09). The CYPD3176 is a 24-pin QFN USB-C sink
controller. It contains the CC terminations and negotiates the power contract.
It drives an external P-channel pass switch when the contract is within the
programmed voltage window. Its D+/D$-$ pins support charging protocols for
sources without USB Power Delivery.

{\footnotesize
\begin{longtable}{|L{4.4cm}|c|c|c|c|L{3.6cm}|c|}
\caption{CYPD3176 electrical characteristics}\label{tab:cypd3176}\\
\hline
\textbf{Parameter} & \textbf{Min} & \textbf{Typ} & \textbf{Max} & \textbf{Units} & \textbf{Conditions} & \textbf{Page} \\
\hline
\endfirsthead
\caption[]{CYPD3176 electrical characteristics (continued)}\\
\hline
\textbf{Parameter} & \textbf{Min} & \textbf{Typ} & \textbf{Max} & \textbf{Units} & \textbf{Conditions} & \textbf{Page} \\
\hline
\endhead
\hline
\endlastfoot
VBUS\_IN supply input voltage & \num{3.0} & & \num{24.5} & \si{\volt} & \(-40\)~\si{\celsius} \(\leq T_A \leq\) \SI{105}{\celsius} & 18 \\
\hline
VDDD supply input voltage & \num{2.7} & & \num{5.5} & \si{\volt} & sink mode; \(-40\)~\si{\celsius} \(\leq T_A \leq\) \SI{105}{\celsius} & 18 \\
\hline
VCCD core logic output & & \num{1.8} & & \si{\volt} & & 18 \\
\hline
Supply current from VBUS or VDDD & & \num{10} & & \si{\milli\ampere} & VDDD = 5~V OR VBUS = 5~V, \(T_A\) = \SI{25}{\celsius}; CC1/CC2 in Tx or Rx, no I/O sourcing current, 2 SCBs at 1~Mbps, EA/ADC/CSA/UVOV ON, CPU at 24~MHz & 18 \\
\hline
VDDD decoupling capacitor & \num{0.8} & \num{1} & & \si{\micro\farad} & X5R ceramic or better & 18 \\
\hline
VBUS\_IN\_DIS decoupling capacitor & \num{0.1} & & & \si{\micro\farad} & X5R ceramic or better & 18 \\
\hline
VBUS\_IN\_DIS, DC\_OUT\_DIS absolute maximum & & & \num{30} & \si{\volt} & absolute max & 16 \\
\hline
VDDD absolute maximum & & & \num{6} & \si{\volt} & & 16 \\
\hline
CC1, CC2 absolute maximum & & & \num{22} & \si{\volt} & as per USB PD specification, maximum allowed VBUS = 21.5~V & 16 \\
\hline
VBUS\_CTRL absolute maximum & \num{-0.3} & & \multicolumn{2}{c|}{VBUS\_MAX} & & 17 \\
\hline
GPIO absolute maximum & \num{-0.5} & & \multicolumn{2}{c|}{VDDD + \num{0.5}} & & 17 \\
\hline
Rd UFP CC termination & \num{4.59} & \num{5.1} & \num{5.61} & \si{\kilo\ohm} & & 22 \\
\hline
Gate-driver pull-down resistance & & & \num{3} & \si{\kilo\ohm} & applicable on VBUS\_CTRL and SAFE\_PWR\_EN to turn ON external PFET & 24 \\
\hline
Gate-driver pull-up resistance & & & \num{4} & \si{\kilo\ohm} & applicable on VBUS\_CTRL to turn OFF external PFET & 24 \\
\hline
Gate-driver pull-down current, strength 1 & \num{25} & & \num{75} & \si{\micro\ampere} & I-mode (current mode) pull down at 5~V & 24 \\
\hline
Gate-driver pull-down current, strength 8 & \num{280} & & \num{580} & \si{\micro\ampere} & & 24 \\
\hline
CSP input capacitance & \num{7} & & \num{10} & \si{\pico\farad} & guaranteed by characterization & 22 \\
\hline
D+, D$-$ pin capacitance & & & \num{22} & \si{\pico\farad} & capacitance on D+, D\(-\) pins; guaranteed by characterization & 19 \\
\hline
ESD, human body model & \num{2200} & & & \si{\volt} & & 16 \\
\hline
ESD, IEC 61000-4-2 contact & \num{8000} & & & \si{\volt} & contact discharge on CC1, CC2, VBUS\_IN\_DIS, HPI\_SDA and HPI\_SCL pins & 16 \\
\hline
ESD, IEC 61000-4-2 air & \num{15000} & & & \si{\volt} & air discharge for D+, D\(-\), CC1, CC2, VBUS\_IN\_DIS, HPI\_SDA and HPI\_SCL pins & 16 \\
\hline
\end{longtable}
}

The datasheet gives this negotiation order
(\ds{infineon-cypd3176-24lqxq-datasheet-en.pdf}{15}):

\begin{quote}
Negotiation sequence: (PD $\rightarrow$ BC1.2 $\rightarrow$ AFC $\rightarrow$
QC2.0 (Class B) $\rightarrow$ Apple $\rightarrow$ Type-C Only) \dots When
negotiating legacy protocols, the BCR-PLUS device negotiates voltage and current
based on the VBUS\_MIN/MAX and ISNK\_COARSE/ISNK\_FINE settings.
\end{quote}

It also specifies the alternate power path
(\ds{infineon-cypd3176-24lqxq-datasheet-en.pdf}{14}):

\begin{quote}
This gate driver is enabled when the PD contract defaults to 5~V due to
mismatching capabilities. This driver is not enabled when the device is
configured so that 5~V input is accepted by the system.
\end{quote}

Accepting a \SI{5}{\volt} input disables the SAFE\_PWR\_EN path. This design
accepts \SI{5}{\volt}, so it does not use SAFE\_PWR\_EN.

{\sloppy
The register interface through which the MCU reaches the part is specified
separately, in
\texttt{infineon-\allowbreak ez-\allowbreak pd-\allowbreak tm-\allowbreak
bcr-\allowbreak plus-\allowbreak bcr-\allowbreak lite-\allowbreak
host-\allowbreak processor-\allowbreak interface-\allowbreak
specification-\allowbreak cypd3176-\allowbreak cypd3178-\allowbreak
usermanual-\allowbreak en.pdf}
(Infineon, 002-33575 Rev.~*A, 2023-10), cited as the HPI specification.
SINK\_RDO\_REQUEST sends a Request Data Object to the source
(\cref{sec:pd-sequence}). CHARGING\_MODE\_STATUS identifies the active charging
protocol (\cref{sec:pd-data}).
\par}

\Cref{sec:pd-data} gives the USB 2.0 analog-switch requirements.

\subsection{Load budget referred to the bus}\label{sec:pd-load}

The maximum input powers are:
\begin{align*}
  P_{\mathrm{in,HV}} &= \SI{3.54}{\watt} && \text{(\cref{sec:hv-op}, design-max envelope, } \eta = 0.833), \\
  P_{\mathrm{in,5V}} &= \SI{0.588}{\watt} && \text{(\cref{sec:5v-ss}, \SI{0.50}{\watt} at } \eta = 0.85), \\
  P_{\mathrm{in,3V3}} &= \SI{1.39}{\watt} && \text{(\cref{sec:3v3-load}, \SI{1.18}{\watt} at } \eta = 0.85),
\end{align*}
The additive maximum is \SI{5.51}{\watt}. At each contract's $-5\,\%$ limit,
include the controller's \SI{10}{\milli\ampere} current:
\begin{align*}
  I_{\mathrm{port}} &= \frac{5.51}{V_{\mathrm{BUS}}} + 0.010
  &&\SI{4.75}{\volt}: \SI{1.170}{\ampere}; \quad
    \SI{8.55}{\volt}: \SI{0.654}{\ampere}; \quad
    \SI{11.4}{\volt}: \SI{0.493}{\ampere}.
\end{align*}
The \SI{4.75}{\volt} case sets the conducted-current requirements. The maximum
condition has all digit tubes at maximum current during continuous Wi-Fi
transmission. This is a design limit, not an expected operating point.

\subsection{Contract window}\label{sec:pd-window}

VBUS\_MIN and VBUS\_MAX each take one of six settings through a divider from
VDDD (\ds{infineon-cypd3176-24lqxq-datasheet-en.pdf}{13}), listed in
\cref{tab:pd-vbus-settings}. The pair defines the range of source voltages the
controller will accept.

\begin{table}[htbp]
\centering
\caption{VBUS\_MIN and VBUS\_MAX settings}\label{tab:pd-vbus-settings}
\begin{tabular}{|c|c|c|c|c|}
\hline
\textbf{Voltage (V)} & \textbf{Ratio to VDDD} & \textbf{Pull-up (\si{\kilo\ohm})} & \textbf{Pull-down (\si{\kilo\ohm})} & \textbf{Pin window (mV)} \\
\hline
5 & 0/6 & open & 0 & 0 to 248 \\
\hline
9 & 1/6 & 5.1 & 1 & 249 to 786 \\
\hline
12 & 2/6 & 5.1 & 2.4 & 787 to 1347 \\
\hline
15 & 3/6 & 5.1 & 5.1 & 1348 to 1920 \\
\hline
19 & 4/6 & 5.1 & 10 & 1921 to 2778 \\
\hline
20 & $\geq$5/6 & 0 & open & $\geq$2779 \\
\hline
\end{tabular}
\end{table}

\paragraph{Upper limit.} The downstream converters cannot accept the
\SI{15}{\volt} setting. A
\SI{15}{\volt} contract at its $+5\,\%$ limit is \SI{15.75}{\volt}, which is
\SI{87.5}{\percent} of the TPS629203's \SI{18}{\volt} absolute maximum
(\cref{sec:tps629203}) and past the \SI{80}{\percent} derating this document
applies to voltage ratings. Set VBUS\_MAX to \SI{12}{\volt}. At $+5\,\%$, this
setting gives \SI{12.6}{\volt}. This value is the converter input maximum used
throughout the analysis.

\paragraph{Lower limit.} The board must start from a USB host that supplies
only \SI{5}{\volt}. Set VBUS\_MIN to \SI{5}{\volt} by connecting the pin to
ground. Set VBUS\_MAX to \SI{12}{\volt} through
\des{R}{PDMAX1} $= \SI{5.1}{\kilo\ohm}$ to VDDD and \des{R}{PDMAX2} $=
\SI{2.4}{\kilo\ohm}$ to ground.

The accepted voltages are \SI{5}{\volt}, \SI{9}{\volt}, and
\SI{12}{\volt}. The controller requests the highest available voltage within
the window (inferred). The datasheet does not specify selection within the
window. Firmware can request another contract if this inference is incorrect
(\cref{sec:pd-sequence}).

\subsection{Operating current request}\label{sec:pd-isnk}

ISNK\_COARSE and ISNK\_FINE set the requested current
(\ds{infineon-cypd3176-24lqxq-datasheet-en.pdf}{13, 14}). The coarse setting
uses \SI{1}{\ampere} steps from \SIrange{0}{5}{\ampere}. The fine settings are
\SI{0}{}, \SI{250}{}, \SI{500}{}, \SI{750}{}, and
\SI{900}{\milli\ampere}. The initial request must permit operation from a
\SI{500}{\milli\ampere} host port.

A fixed request cannot support full display operation and a
\SI{500}{\milli\ampere} host port. Full operation requires at least
\SI{750}{\milli\ampere} at \SI{9}{\volt}. It requires more than
\SI{900}{\milli\ampere} at \SI{5}{\volt} (\cref{tab:pd-required}). A higher
initial request would prevent a \SI{500}{\milli\ampere} host contract.

The resistor settings request the current for limited operation. Firmware
requests full-operation current after start-up (\cref{sec:pd-sequence}).
Limited operation draws \SI{0.404}{\ampere} (\cref{sec:pd-classes}). The
\SI{500}{\milli\ampere} setting gives \SI{24}{\percent} margin. Set
ISNK\_COARSE to
\SI{0}{\ampere}, tied to ground; ISNK\_FINE = $+\SI{500}{\milli\ampere}$
through \des{R}{PDI1} $= \SI{5.1}{\kilo\ohm}$ to VDDD and \des{R}{PDI2} $=
\SI{2.4}{\kilo\ohm}$ to ground.

The initial request permits \SI{2.5}{\watt}, \SI{4.5}{\watt}, or
\SI{6}{\watt}. The display inhibit prevents full display operation before
firmware confirms sufficient current.

\subsection{Operating classes}\label{sec:pd-classes}

All converters can operate from each accepted voltage. The available current
determines the operating class. Full operation requires the current in
\cref{tab:pd-required}.

\begin{table}[htbp]
\centering
\caption{Full-operation current requirement and firmware request, by
contract}\label{tab:pd-required}
\begin{tabular}{|c|c|c|c|}
\hline
\textbf{Contract} & \textbf{Draw at $-5\,\%$ (\cref{sec:pd-load})} &
\textbf{Request} & \textbf{Margin} \\
\hline
\SI{5}{\volt} & \SI{1.170}{\ampere} & \SI{1.25}{\ampere} & $7\,\%$ \\
\hline
\SI{9}{\volt} & \SI{0.654}{\ampere} & \SI{750}{\milli\ampere} & $15\,\%$ \\
\hline
\SI{12}{\volt} & \SI{0.493}{\ampere} & \SI{750}{\milli\ampere} & $52\,\%$ \\
\hline
\end{tabular}
\end{table}

\paragraph{Full operation.} The source must guarantee the request in
\cref{tab:pd-required}. A PD contract gives this guarantee directly. A BC1.2
charging port supplies at least \SI{1.5}{\ampere}. QC~2.0 and AFC sources
supply the required current at \SI{9}{\volt} or \SI{12}{\volt} (inferred).
A USB-C source that advertises \SI{1.5}{\ampere} or \SI{3}{\ampere} supplies
the \SI{5}{\volt} requirement. Firmware then releases the display inhibit.

\paragraph{Limited operation.} Limited operation applies when the source does
not guarantee the table current. The display inhibit remains asserted. The
flyback remains off. The drivers and clock controller remain operational.

The load in limited operation is the 3.3~V rail and the idle drivers:
\begin{align*}
  I_{\mathrm{port}}^{\mathrm{lim}} &= \frac{1.39}{4.72} + 0.100 + 0.010
    = \SI{0.404}{\ampere},
\end{align*}
The terms are the 3.3~V converter, the drivers, and the controller. This is
\SI{81}{\percent} of a \SI{500}{\milli\ampere} host-port supply.

\subsection{Start-up and renegotiation}\label{sec:pd-sequence}

The initial contract supplies limited operation. Firmware must confirm a
higher-current contract before it enables the display.

\begin{enumerate}[itemsep=2pt]
  \item \textbf{Attach.} The controller runs from the receptacle's VBUS through
    its own regulator and negotiates the highest voltage in the window at
    \SI{500}{\milli\ampere}. VBUS\_CTRL closes under the slew control of
    \cref{sec:pd-fet}.
  \item \textbf{Rails.} The 3.3~V rail comes up and the MCU boots; the 5~V rail
    comes up and the drivers idle. The flyback stays off behind the inhibit of
    \cref{sec:hv-inhibit}, which no contract voltage bypasses.
  \item \textbf{Interrogate.} The MCU reads the negotiated result, the active
    legacy protocol, and the measured bus voltage over HPI
    (\cref{sec:pd-pins}).
  \item \textbf{Renegotiate.} On a PD contract short of full operation, the MCU
    requests the operating current \cref{tab:pd-required} assigns that
    contract's voltage, covering the referred draw with \SI{7}{\percent} to
    spare at the tightest point.
  \item \textbf{Release.} On a result that guarantees its voltage's requirement
    (\cref{sec:pd-classes}) the MCU releases the display inhibit. On any other
    outcome it does not, and the board stays in limited operation for as long
    as that source is attached.
\end{enumerate}

\paragraph{Display inhibit.} The flyback UVLO releases at
\SI{4.4}{\volt} (\cref{sec:hv-uvlo}). Every accepted contract exceeds this
voltage. Therefore, UVLO cannot identify a sufficient-current contract. The
MCU controls a default-active display inhibit (\cref{sec:hv-inhibit}).

\paragraph{Current-only renegotiation.} If the controller selected the highest
available voltage, step 4 changes only the current. The pass switch remains
closed. A voltage change can charge the \SI{135}{\micro\farad} bus capacitance
without pass-switch slew control. The source can current-limit or return to its
safe voltage.

If the controller does not select the highest voltage, step 4 can cause one
voltage transition. The hardware must tolerate this event.

\paragraph{Step 4 request.} A Request Data Object identifies the selected
source capability and operating current. SINK\_RDO\_REQUEST sends the 32-bit
object to the source: ``Writing to this
register will result in Request of PD contract re-negotiation'' (HPI
specification, p.~24). Do not replace the sink capability list through
SELECT\_SINK\_PDO. The specification warns that ``once the Sink PDOs are
updated, BCR-PLUS will re-negotiate the PD
contract with the power adapter. This will cause momentary power loss to the
system'' (HPI specification, p.~22). Step 4 writes an RDO and retains the
resistor-defined capability list.

The sink PDO mask is nonzero at reset. A SUCCESS response confirms only that
the controller queued the request. Step 5 waits for the
contract-negotiation-complete response and checks its accepted request.

\subsection{Bus pass switch}\label{sec:pd-fet}

VBUS\_CTRL drives \des{Q}{BUS}'s gates through \des{R}{BUSG}, asserting only when
a contract inside the window closes.

\paragraph{Bidirectional blocking.} A single P-channel MOSFET blocks current in
only one direction when off. Its body diode can discharge the
\SI{135}{\micro\farad} bus capacitance into the receptacle after detach. It can
also conduct a receptacle overvoltage into $V_{\mathrm{BUS}}$.

Use two P-channel MOSFETs with common sources and common gates. Their opposing
body diodes block current in both directions. Connect DC\_OUT\_DIS to the
switch output (\ds{infineon-cypd3176-24lqxq-datasheet-en.pdf}{9}).

\paragraph{Ratings.} Each MOSFET must withstand a \SI{24}{\volt} ESD clamp
voltage (\cref{sec:pd-esd}). With \SI{20}{\percent} derating, require
$V_{DS}\geq\SI{30}{\volt}$. The maximum continuous current is
\SI{1.170}{\ampere}. Require $I_D\geq\SI{1.5}{\ampere}$.

\paragraph{On-resistance.} Require
$R_{DS(\mathrm{on})}\leq\SI{25}{\milli\ohm}$ per device at the available gate
drive. The pair drops \SI{59}{\milli\volt} at \SI{1.170}{\ampere} and
dissipates \SI{68}{\milli\watt}.

The \SI{4.75}{\volt} converter input floor applies at $V_{\mathrm{BUS}}$. A
\SI{5}{\volt} contract at its $-5\,\%$ limit has no voltage-drop budget. A
\SI{250}{\milli\ohm} cable drops \SI{293}{\milli\volt} at
\SI{1.170}{\ampere}. With the pass-switch drop, $V_{\mathrm{BUS}}$ can fall to
approximately \SI{4.40}{\volt}. The flyback UVLO falling threshold is
\SI{3.97}{\volt} (\cref{sec:hv-uvlo}). If the input falls below this threshold,
the display enters limited operation.

\paragraph{Inrush.} The three converters provide \SI{135}{\micro\farad} of
nominal input capacitance: \SI{100}{\micro\farad} polymer and
\SI{4.7}{\micro\farad} ceramic at the flyback (\cref{sec:hv-caps}),
\SI{10}{\micro\farad} at the 5~V converter (\cref{sec:5v-cin}), and
\SI{20}{\micro\farad} across the 3.3~V converter's two input capacitors
(\cref{sec:3v3-cin}). The attach contract supplies
\SI{500}{\milli\ampere}. Limited operation uses \SI{404}{\milli\ampere}.
Thus, \SI{96}{\milli\ampere} remains for inrush. A $+20\,\%$ capacitance
tolerance gives \SI{162}{\micro\farad}. Use this value for the inrush limit:
\begin{align*}
  t_{\mathrm{ramp}} &\geq \frac{C\,\Delta V}{I_{\mathrm{inrush}}}
    = \frac{\SI{162}{\micro\farad}\times\SI{12.6}{\volt}}{\SI{0.096}{\ampere}}
    = \SI{21.3}{\milli\second}.
\end{align*}
Require $t_{\mathrm{ramp}}\geq\SI{22}{\milli\second}$. The gate network sets
this ramp.

\paragraph{Gate network.} Five components hold the switch off before controller
start-up. They also limit gate-source voltage and set the output ramp
(\ds{infineon-cypd3176-24lqxq-datasheet-en.pdf}{12}).

\begin{description}
  \item[\des{R}{BUSOFF}, gate to common source] holds $V_{GS}$ at zero when the
    controller does not drive VBUS\_CTRL. Set it to \SI{49.9}{\kilo\ohm}. It
    carries \SI{200}{\micro\ampere} at full gate drive.
  \item[\des{R}{BUSBIAS}, gate to the receptacle's VBUS] biases the switch off
    during connection. With \des{R}{BUSOFF}, it sets the off-state $V_{GS}$ to
    approximately $+\SI{0.6}{\volt}$. Set it to \SI{10}{\kilo\ohm}.
  \item[\des{D}{BUSG}, cathode on the source, anode on the gate] bounds
    $V_{GS}$. It is a protective element: with the driver's
    \SI{3}{\kilo\ohm} pull-down working against \des{R}{BUSBIAS} and
    \des{R}{BUSOFF}, the gate settles at
    \begin{align*}
      V_G &= V_S\,\frac{G_{\mathrm{BIAS}}+G_{\mathrm{OFF}}}
                       {G_{\mathrm{BIAS}}+G_{\mathrm{OFF}}+G_{\mathrm{PD}}}
        = 12.6 \times 0.265 = \SI{3.3}{\volt},
    \end{align*}
    This gives $V_{GS}=-\SI{9.3}{\volt}$ during normal operation. Use a
    \SI{10}{\volt} Zener clamp. With \SI{20}{\percent} derating, require
    $|V_{GS}|\geq\SI{12.5}{\volt}$ for \des{Q}{BUS}.
  \item[\des{C}{BUSSLEW}, gate to $V_{\mathrm{BUS}}$] sets the output ramp. It
    sets $dV_{\mathrm{BUS}}/dt=I_{PD}/C$. Driver-strength setting 1 sinks
    \SIrange{25}{75}{\micro\ampere}. The required effective capacitance is
    \begin{align*}
      C_{\mathrm{BUSSLEW}} &\geq
        \frac{I_{PD}^{\max}t_{\mathrm{ramp}}^{\min}}{\Delta V}
        = \frac{\SI{75}{\micro\ampere}\times\SI{21.3}{\milli\second}}
               {\SI{12.6}{\volt}}
        = \SI{127}{\nano\farad}.
    \end{align*}
    Require \SI{130}{\nano\farad} effective at \SI{12.6}{\volt}. The
    installed \SI{1}{\micro\farad} 0402 part is expected to provide
    approximately \SI{200}{\nano\farad} under bias, as stated on the schematic.
    At \SI{200}{\nano\farad}, the slew rate is
    \SIrange{0.125}{0.375}{\volt\per\milli\second}. The \SI{12.6}{\volt} ramp
    takes \SIrange{33.6}{100.8}{\milli\second}. The maximum inrush is
    \SI{61}{\milli\ampere} with the \SI{162}{\micro\farad} tolerance stack.
    Limited-operation current plus inrush is \SI{465}{\milli\ampere}.
  \item[\des{C}{BUSCM}, gate to the receptacle's VBUS] holds $V_{GS}$ during a
    fast input transient. Its installed value is \SI{1}{\micro\farad}. The
    schematic expects approximately \SI{200}{\nano\farad} under bias.
\end{description}

\des{R}{BUSG} connects VBUS\_CTRL to the gate node. It limits clamp current.
Set \des{R}{BUSG} to \SI{1}{\kilo\ohm}.

The drive-strength setting controls the ramp. Configure setting 1. Its full
current range satisfies the inrush limit with the specified effective
\des{C}{BUSSLEW} value.

\subsection{Bus-rail discharge}\label{sec:pd-bleed}

\Cref{sec:3v3-diodeor} requires
$R_{\mathrm{BUSBLEED}}\leq\SI{287}{\kilo\ohm}$. The controller discharge FET
does not satisfy this continuous requirement.

The controller contains a VBUS discharge FET, described as
``used to discharge VBUS to meet the USB-PD specification timing on a detach
condition'' (\ds{infineon-cypd3176-24lqxq-datasheet-en.pdf}{5}). The datasheet
does not specify a static impedance. During header-powered operation,
\des{Q}{3V3BUS} leakage continuously charges the bus rail.

Set \des{R}{BUSBLEED} to \SI{100}{\kilo\ohm}. At \SI{12.6}{\volt}, it draws
\SI{126}{\micro\ampere} and dissipates \SI{1.59}{\milli\watt}.

\subsection{Data path}\label{sec:pd-data}

D+ and D$-$ connect to the controller, the ESD array, and the MCU USB interface
(\cref{sec:esp32-usb}). An analog switch isolates the MCU branch during
charging-protocol negotiation.

\paragraph{Charging-protocol negotiation.} QC~2.0 signals its
\SI{9}{\volt} and \SI{12}{\volt} levels on D+ and D$-$. These levels match the
accepted voltage window. Therefore, CHARGING\_MODE enables all supported
charging protocols (\cref{sec:pd-pins}).

\paragraph{Data-pair ownership.} QC~2.0 and AFC hold their contract state on
D+ and D$-$. The controller must retain the pair while these protocols are
active. The MCU cannot use USB data in this condition.

USB hosts use USB Power Delivery or USB-C current advertisement (inferred).
Dedicated chargers can use QC~2.0 or AFC. Firmware identifies the active mode
before it connects the MCU branch.

CHARGING\_MODE\_STATUS reports no charging protocol, BC1.2, QC~2.0, AFC, or
Apple mode (HPI specification, p.~18). Firmware reads this register before it
closes the switch. It closes the switch for no charging protocol or BC1.2. It
keeps the switch open for QC~2.0, AFC, or Apple mode.

\paragraph{MCU-branch switch.} The controller requires D+ and D$-$ before the
3.3~V rail starts. Therefore, a 3.3~V analog switch must not isolate the
controller. Place a single-throw switch only in the MCU branch. Its unpowered
state is open. Firmware closes it before USB enumeration.

The MCU USB pins have clamps to the 3.3~V rail. An unswitched MCU branch can
load D+ and D$-$ while this rail is off. Both pins also pulse high during MCU
start-up (\cref{sec:esp32-usb}). The open switch isolates these pulses.

\des{R}{USBSW} holds the switch disabled while the MCU pin is high impedance.
Use an MCU pin without a strapping function (\cref{sec:esp32-usb}). Tie the
switch-select input to the populated branch.

\paragraph{Orientation.} The USB~2.0 receptacle connects A6/A7 to B6/B7.
Therefore, the data path does not require orientation switching. The design
does not use the FLIP output.

\paragraph{Data-line capacitance.} The USB~2.0 upstream-port limit is
\SI{100}{\pico\farad} per line. \Cref{tab:pd-datacap} gives the maximum
contributors.

\begin{table}[htbp]
\centering
\caption{Capacitance on one line of the data pair}\label{tab:pd-datacap}
\begin{tabular}{|l|r|l|}
\hline
\textbf{Contributor} & \textbf{pF} & \textbf{Source} \\
\hline
CYPD3176 D+ or D$-$ pin & \num{22} & \cref{tab:cypd3176} \\
\hline
\des{D}{USBESD} & \num{3.5} & \cref{sec:pd-esd} \\
\hline
Switch common side, closed & \num{7.5} & \cref{tab:pd-switch} \\
\hline
MCU transceiver and traces & \num{15} & estimate \\
\hline
\textbf{Total} & \textbf{\num{48}} & against \SI{100}{\pico\farad} \\
\hline
\end{tabular}
\end{table}

The MCU and trace value is an estimate because the module datasheet does not
specify USB-pin capacitance. The total has \SI{52}{\pico\farad} of margin.

This result applies only to USB 2.0 full-speed operation. The controller input
capacitance exceeds a typical high-speed budget (inferred).

\paragraph{Switch requirements.} The branch carries USB~2.0 full speed at
\SI{12}{\mega\bit\per\second}. The switch must remain open while unpowered. Its
I/O pins must not clamp the data lines to its unpowered supply.

\begin{table}[htbp]
\centering
\caption{Data-path switch requirements}\label{tab:pd-switch}
\begin{tabular}{|l|L{8.6cm}|}
\hline
\textbf{Parameter} & \textbf{Requirement} \\
\hline
Rating & USB 2.0 full speed, \SI{12}{\mega\bit\per\second} \\
\hline
Bandwidth & $\geq\SI{50}{\mega\hertz}$ at \num{-3}~dB, for
  \SIrange{4}{20}{\nano\second} edges \\
\hline
On-capacitance & $\leq\SI{15}{\pico\farad}$ \\
\hline
On-resistance & $\leq\SI{10}{\ohm}$, channel mismatch $\leq\SI{1}{\ohm}$;
  it adds to the transceiver's \SIrange{28}{44}{\ohm} source impedance \\
\hline
Signal range & $\geq\SI{3.6}{\volt}$, covering the legacy negotiation levels
  the open switch is exposed to \\
\hline
Supply & \SI{3.3}{\volt}, from the MCU's own rail \\
\hline
Unpowered state & open, and drawing no more than a few \si{\micro\ampere} with
  its supply at \SI{0}{\volt} and either I/O anywhere up to
  \SI{3.6}{\volt} \\
\hline
\end{tabular}
\end{table}

The TS3USB221A (\texttt{ts3usb221a.pdf}, Texas Instruments, SCDS277C,
2024-10) meets these requirements. It has \SI{900}{\mega\hertz} typical bandwidth,
\SI{7.5}{\pico\farad} maximum on-capacitance, \SI{6}{\ohm} maximum on-resistance
with \SI{0.2}{\ohm} typical mismatch, signal handling to \SI{5.5}{\volt},
\SIrange{2.3}{3.6}{\volt} supply, and $\pm\SI{2}{\micro\ampere}$ power-off
leakage with $V_{CC} = \SI{0}{\volt}$ and its I/O anywhere from \num{0} to
\SI{5.25}{\volt}. Use it as a single-throw switch with the second branch
unconnected.

Its enable input disconnects on a logic high. Therefore, \des{R}{USBSW} pulls
the input up to the 3.3~V rail. The MCU drives the input low to connect USB.

\paragraph{Controller branch.} \des{R}{PDDP} and \des{R}{PDDM} connect the
controller to D+ and D$-$. Set both to \SI{0}{\ohm}.

The controller releases D+ and D$-$ when no charging protocol is active
(inferred). The datasheet does not specify the idle state. If testing shows an
active drive, remove \des{R}{PDDP} and \des{R}{PDDM}. This change disables
charging protocols that require D+ and D$-$, but it retains USB data and USB
Power Delivery.

\subsection{ESD protection}\label{sec:pd-esd}

The receptacle requires IEC 61000-4-2 protection. D+ and D$-$ operate at logic
levels. CC operates near ground. VBUS can reach \SI{12.6}{\volt}. Separate
devices protect these groups.

The controller specifies \SI{8}{\kilo\volt} contact discharge on CC1, CC2, and
VBUS\_IN\_DIS (\ds{infineon-cypd3176-24lqxq-datasheet-en.pdf}{16}). It does not
specify contact discharge on D+ or D$-$.

\paragraph{CC1 and CC2.} Do not add external clamps. Only the rated controller
pins connect to these lines.

\paragraph{D+ and D$-$.} Use a matched ESD array. It must meet IEC 61000-4-2
level 4. Its capacitance must not exceed \SI{4}{\pico\farad} per line.

\des{D}{USBESD}, the USBLC6-2 (\texttt{USBLC6-2.pdf}, STMicroelectronics),
satisfies these requirements. It has \SI{3.5}{\pico\farad} maximum line capacitance,
\SI{0.015}{\pico\farad} matching between the two I/O lines,
\SI{150}{\nano\ampere} maximum leakage, and \SI{8}{\kilo\volt} contact with
\SI{15}{\kilo\volt} air discharge, in a SOT-666 or SOT23-6L body.

Connect its rail pin to VDDD. The pin has a \SI{5.25}{\volt} standoff rating
and a \SI{6}{\volt} minimum breakdown voltage (\ds{USBLC6-2.pdf}{2}). Therefore,
it cannot connect to a \SI{9}{\volt} or \SI{12}{\volt} VBUS.

\paragraph{VBUS.} Use a separate transient suppressor. It must have a standoff
above \SI{12.6}{\volt} and meet IEC 61000-4-2 level 4. Its clamp voltage must
remain below all connected-device absolute maximum ratings.

Use a unidirectional device. VBUS\_IN\_DIS has a
$-\SI{0.3}{\volt}$ absolute minimum. A unidirectional device limits a negative
event to its forward voltage.

\des{D}{BUSTVS}, the SD15Q (\texttt{SD15Q.pdf}, Diodes Incorporated), meets these requirements:
\SI{15}{\volt} standoff against \SI{12.6}{\volt} of operating voltage,
\SI{16}{\volt} minimum breakdown, \SI{1}{\micro\ampere} maximum leakage at the
full standoff, and $\pm\SI{30}{\kilo\volt}$ contact and air discharge to
IEC 61000-4-2, in SOD-323.

Use the \SI{100}{\nano\second} transmission-line-pulse curve for the ESD event
(\ds{SD15Q.pdf}{3}). It gives approximately \SI{19.5}{\volt} at turn-on and
\SI{24}{\volt} at \SI{30}{\ampere}. The dynamic resistance is approximately
\SI{0.15}{\ohm} (inferred).

\paragraph{Clamp-voltage effects.} All components on the node must withstand
\SI{24}{\volt} during an ESD event. Apply absolute maximum ratings to this
transient condition.
\begin{description}
  \item[\des{Q}{BUS}] Require $V_{DS}\geq\SI{30}{\volt}$. The device must also
    meet the \SI{25}{\milli\ohm} resistance requirement.
  \item[\des{C}{PDIN}] sees the clamp across it. A \SI{25}{\volt} part is
    under-rated for a \SI{24}{\volt} transient. Use a \SI{50}{\volt} part.
  \item[VBUS\_IN\_DIS] has a \SI{30}{\volt} absolute maximum, used to
    \SI{80}{\percent} at the clamp. The pin also has an integrated
    \SI{8}{\kilo\volt} contact-discharge rating.
\end{description}

\subsection{Configuration and housekeeping pins}\label{sec:pd-pins}

\paragraph{CHARGING\_MODE.} Leave this pin open to select all charging
protocols. This permits QC~2.0 negotiation at \SI{9}{\volt} and
\SI{12}{\volt}. A \SI{5}{\kilo\ohm} pull-up would select only BC1.2.

The controller samples this strap at power-up
(\ds{infineon-cy4534-bcr-plus-usermanual-en.pdf}{14}). Firmware does not write
CHARGING\_\allowbreak MODE\_\allowbreak DISABLE.

\paragraph{DATAMODE/FLIP.} Leave this pin open. This is the configuration for
``applications that need to sink power as well as support USB communication''
(\ds{infineon-cypd3176-24lqxq-datasheet-en.pdf}{9}). A
\SI{1}{\kilo\ohm} pull-up selects power-only operation.

Firmware controls data-pair ownership as specified in \cref{sec:pd-data}.

\paragraph{SAFE\_PWR\_EN.} Leave this pin open. The accepted \SI{5}{\volt}
window disables this function.

\paragraph{FAULT.} Connect this pin to the MCU. Set
\des{R}{PDFLT}=\SI{100}{\kilo\ohm} to ground. This selects the
\SI{3.3}{\volt} push-pull mode
(\ds{infineon-cypd3176-24lqxq-datasheet-en.pdf}{8}). The signal reports
overvoltage, overcurrent, and overtemperature faults without an HPI
transaction.

\paragraph{Fault indicator.} The push-pull output can also drive an indicator.
The pin guarantees
$V_{OH} \geq V_{DDD} - \SI{0.6}{\volt}$ at \SI{4}{\milli\ampere}
(\ds{infineon-cypd3176-24lqxq-datasheet-en.pdf}{19}), an output resistance of
\SI{150}{\ohm}. The MCU high threshold is \SI{2.475}{\volt}
(\cref{sec:esp32-mcu}). The indicator current is:
\begin{align*}
  I &= \frac{V_{DDD} - V_F}{R + \SI{150}{\ohm}}, &
  V_{OH} &= V_{DDD} - \SI{150}{\ohm}\cdot I,
\end{align*}
which at $R = \SI{1}{\kilo\ohm}$ and a red diode's \SI{1.9}{\volt} gives
\SI{1.22}{\milli\ampere} and \SI{3.12}{\volt}. This gives
\SI{0.65}{\volt} of input margin. Require $\des{R}{PDIND}\geq
\SI{1}{\kilo\ohm}$. Use a red or amber indicator LED.

\paragraph{HPI.} Connect HPI\_SDA, HPI\_SCL, and HPI\_INT to the MCU. The
controller is an I\textsuperscript{2}C slave at address 0x08. Pull each
open-drain signal up to the 3.3~V rail with \SI{10}{\kilo\ohm}. Firmware uses
HPI to request full-operation current.

\subparagraph{Pull-up supply.} VDDD starts before the MCU 3.3~V rail. VDDD
pull-ups would drive unpowered MCU inputs. Therefore, connect the pull-ups to
the MCU 3.3~V rail.

\subparagraph{Pull-up value.} Fast mode permits a \SI{300}{\nano\second} rise
time. Use \SI{27}{\pico\farad} of total bus capacitance (inferred). A
\SI{10}{\kilo\ohm} pull-up gives \SI{229}{\nano\second}. This supports the
controller's \SI{400}{\kilo\bit\per\second} maximum rate.

\paragraph{Supply and decoupling.} Connect VBUS\_IN\_DIS to receptacle VBUS.
Use \des{C}{PDIN}=\SI{100}{\nano\farad}. Use
\des{C}{PDD1}=\SI{1}{\micro\farad} and
\des{C}{PDD2}=\des{C}{PDD3}=\SI{100}{\nano\farad} on VDDD. Use
\des{C}{PDC}=\SI{1}{\micro\farad} on VCCD. The two divider pull-ups draw
\SI{0.98}{\milli\ampere} from VDDD.

\paragraph{CC1 and CC2.} Connect these pins directly to the receptacle. The
controller contains \SIrange{4.59}{5.61}{\kilo\ohm} sink terminations. Do not
add external Rd resistors.

\paragraph{DC\_OUT\_DIS.} Connect this pin to $V_{\mathrm{BUS}}$. Its
\SI{30}{\volt} absolute maximum exceeds the \SI{12.6}{\volt} operating maximum.

\paragraph{DNU1 and DNU2.} Leave these pins open.

\paragraph{CSP.} Connect CSP to VSS. This disables the controller overcurrent
measurement. The reference circuit requires a \SI{5}{\milli\ohm} ground-return
shunt (\ds{infineon-cypd3176-24lqxq-datasheet-en.pdf}{11}). This shunt would
split the receptacle and board ground paths. The negotiated contract and each
converter current limit provide the required current protection. The controller
trip threshold would also require configuration-table programming.

\subsection{Summary}\label{sec:pd-summary}

The analysis uses two inferred controller behaviors. The controller selects the
highest voltage in the programmed window (\cref{sec:pd-window}). It releases D+
and D$-$ when no charging protocol is active (\cref{sec:pd-data}). Hardware
tests must verify both behaviors. A worst-case source, cable, and load can also
reduce $V_{\mathrm{BUS}}$ below \SI{4.75}{\volt} (\cref{sec:pd-fet}). The
flyback UVLO then disables the display.

{\footnotesize
\begin{longtable}{|p{0.23\textwidth}|p{0.335\textwidth}|p{0.335\textwidth}|}
\caption{USB-PD front end: specification summary}\label{tab:pd-summary}\\
\hline
\textbf{Parameter} & \textbf{Derived requirement} & \textbf{Specification / result} \\
\hline
\endfirsthead
\caption[]{USB-PD front end: specification summary (continued)}\\
\hline
\textbf{Parameter} & \textbf{Derived requirement} & \textbf{Specification / result} \\
\hline
\endhead
\hline
\endlastfoot
Bus load & Not applicable & \SI{5.51}{\watt} worst-case additive;
  \SI{1.170}{\ampere} at \SI{4.75}{\volt} \\
\hline
VBUS\_MAX & $\leq\SI{12.6}{\volt}$ at the converters; \SI{80}{\percent} of the
  TPS629203's \SI{18}{\volt} rating & \SI{12}{\volt} setting; \SI{12.6}{\volt}
  at $+5\,\%$; \SI{5.1}{\kilo\ohm}/\SI{2.4}{\kilo\ohm} \\
\hline
VBUS\_MIN & board must start on a source with no high-voltage contract &
  \SI{5}{\volt} setting; pin to ground, no components \\
\hline
Accepted contracts & Not applicable & \SI{5}{\volt}, \SI{9}{\volt}, \SI{12}{\volt} \\
\hline
Operating current, resistors & satisfiable by a \SI{5}{\volt} /
  \SI{500}{\milli\ampere} host port; $\geq\SI{0.404}{\ampere}$ &
  coarse \SI{0}{\ampere} + fine \SI{500}{\milli\ampere};
  \SI{5.1}{\kilo\ohm}/\SI{2.4}{\kilo\ohm} \\
\hline
Operating current, run time & per-contract draw of \cref{tab:pd-required} &
  \SI{1.25}{\ampere} at \SI{5}{\volt}, \SI{750}{\milli\ampere} at \SI{9}{\volt}
  or \SI{12}{\volt}, requested over HPI before the display is released;
  SINK\_RDO\_REQUEST, not SELECT\_SINK\_PDO \\
\hline
Full operation & a result guaranteeing \cref{tab:pd-required}'s current at its
  voltage & any accepted voltage; all three converters in envelope; inhibit
  released \\
\hline
Limited operation & 3.3~V rail in regulation; no display &
  any result short of \cref{tab:pd-required}, including every contract before
  renegotiation; flyback inhibited; \SI{0.404}{\ampere} draw \\
\hline
Display inhibit & default asserted; released only on a confirmed contract &
  MCU shunt across the flyback's UVLO divider (\cref{sec:hv-inhibit}) \\
\hline
\des{Q}{BUS} topology & block in both directions: no bus-capacitance drain onto
  a detached receptacle, no port overvoltage past an open switch & two
  P-channel devices, sources joined, gates driven together \\
\hline
\des{Q}{BUS} ratings & \des{D}{BUSTVS}'s \SI{24}{\volt} clamp on $V_{DS}$;
  \des{D}{BUSG}'s \SI{10}{\volt} clamp on $V_{GS}$; \SI{1.170}{\ampere} &
  $V_{DS}\geq\SI{30}{\volt}$; $|V_{GS}|\geq\SI{12.5}{\volt}$;
  $I_D\geq\SI{1.5}{\ampere}$ \\
\hline
\des{Q}{BUS} on-resistance & drop inside the input-window budget, shared by two
  channels in series & $\leq\SI{25}{\milli\ohm}$ each;
  \SI{59}{\milli\volt}, \SI{68}{\milli\watt} across the pair \\
\hline
Inrush & attach draw plus inrush inside the \SI{500}{\milli\ampere} attach
  contract & \SI{162}{\micro\farad} worst-case bus capacitance;
  \des{C}{BUSSLEW} $\geq\SI{130}{\nano\farad}$ effective; setting 1;
  expected \SIrange{33.6}{100.8}{\milli\second} ramp and
  $\leq\SI{61}{\milli\ampere}$ inrush \\
\hline
Gate network & off before the controller is powered; $V_{GS}$ bounded; ramp
  set deliberately & \des{R}{BUSOFF}/\des{R}{BUSBIAS} bias off at
  $+\SI{0.6}{\volt}$; \des{D}{BUSG} caps $|V_{GS}|$; \des{C}{BUSSLEW} sets
  $dV/dt = I_{PD}/C$ \\
\hline
\des{R}{BUSBLEED} & $\leq\SI{0.66}{\volt}/I_R$ (\cref{sec:3v3-diodeor}) &
  \SI{100}{\kilo\ohm}; the controller's discharge FET is detach-timed, not
  continuous \\
\hline
Data path & one host and one device on the pair; controller holds it at attach &
  single-throw switch in the MCU branch; open unpowered \\
\hline
Data-path arbitration & the MCU must not take the pair while a legacy contract
  is held on it & CHARGING\_MODE\_STATUS read before the switch closes; closed on
  no-legacy-mode or BC1.2 only \\
\hline
Switch & \SI{12}{\mega\bit\per\second}; $\leq\SI{15}{\pico\farad}$;
  $\leq\SI{10}{\ohm}$; isolating at $V_{CC} = \SI{0}{\volt}$ &
  TS3USB221A: \SI{900}{\mega\hertz}, \SI{7.5}{\pico\farad}, \SI{6}{\ohm},
  $\pm\SI{2}{\micro\ampere}$ power-off leakage \\
\hline
Data-pair capacitance & $\leq\SI{100}{\pico\farad}$ per line, USB 2.0 upstream
  port & \SI{48}{\pico\farad} worst-case additive \\
\hline
ESD, D+/D$-$ & IEC 61000-4-2 level 4; $\leq\SI{4}{\pico\farad}$ per line,
  matched & USBLC6-2: \SI{3.5}{\pico\farad}, \SI{0.015}{\pico\farad} matching;
  rail pin on VDDD, not the receptacle \\
\hline
ESD, VBUS & standoff $>\SI{12.6}{\volt}$; unidirectional; clamp survivable by
  the node & SD15Q: \SI{15}{\volt} standoff, \SI{24}{\volt} at the
  \SI{30}{\ampere} first peak, $\pm\SI{30}{\kilo\volt}$ \\
\hline
ESD, CC1/CC2 & Not applicable & none; the controller specifies \SI{8}{\kilo\volt} contact
  on both \\
\hline
CHARGING\_MODE & QC 2.0 reaches \SI{9}{\volt} and \SI{12}{\volt} &
  floating: all legacy protocols; sampled once at power-up \\
\hline
CSP & Not applicable & tied to VSS; sink over-current protection forgone, its shunt
  requiring a split ground plane \\
\hline
DATAMODE/FLIP & data communication required & unconnected: data mode \\
\hline
SAFE\_PWR\_EN & Not applicable & unused; inoperative with \SI{5}{\volt} in the window \\
\hline
FAULT & MCU-readable at \SI{3.3}{\volt} levels & \SI{100}{\kilo\ohm} to ground:
  strong-drive mode \\
\hline
\end{longtable}
}

A \SI{5}{\volt} VBUS\_MIN setting permits start-up from a USB host port. The
setting connects VBUS\_MIN to ground. Firmware enables the display only after
it confirms sufficient contract current.

\subsection{Component selection}\label{sec:pd-parts}

\Cref{tab:pd-parts} gives the component requirements. The ratings are minimum
values. Use the smallest package that satisfies the electrical and thermal
requirements.

{\footnotesize
\begin{longtable}{|L{0.10\textwidth}|L{0.14\textwidth}|L{0.11\textwidth}|L{0.33\textwidth}|L{0.12\textwidth}|}
\caption{USB-PD front end: component selection}\label{tab:pd-parts}\\
\hline
\textbf{Designator} & \textbf{Class} & \textbf{Value} & \textbf{Required ratings and characteristics} & \textbf{Body} \\
\hline
\endfirsthead
\caption[]{USB-PD front end: component selection (continued)}\\
\hline
\textbf{Designator} & \textbf{Class} & \textbf{Value} & \textbf{Required ratings and characteristics} & \textbf{Body} \\
\hline
\endhead
\hline
\endlastfoot
Not applicable & Type-C port controller & CYPD3176 & sink mode; internal Rd; external pass-FET gate driver with current-mode slew control & QFN-24 \\
\hline
\des{J}{USB} & USB Type-C receptacle & To be selected & USB 2.0 pinout with A6/A7 tied to B6/B7; VBUS rated $\geq\SI{16}{\volt}$ and $\geq\SI{1.5}{\ampere}$; through-hole shell anchors preferred for connector retention & Through-hole shell anchors \\
\hline
\des{Q}{BUS} & P-channel MOSFET pair & To be selected & two devices, sources joined and gates common; $V_{DS}\geq\SI{30}{\volt}$, set by \des{D}{BUSTVS}'s clamp; $|V_{GS}|\geq\SI{12.5}{\volt}$, set by \des{D}{BUSG}; $I_D\geq\SI{1.5}{\ampere}$; $R_{DS(\mathrm{on})}\leq\SI{25}{\milli\ohm}$ each; one drain on the receptacle's VBUS, the other on $V_{\mathrm{BUS}}$ & common-source dual, SOT-363 or SOT-563 \\
\hline
Not applicable & USB 2.0 analog switch & TS3USB221A & \cref{tab:pd-switch}; power-off leakage sets selection; open and drawing no more than a few \si{\micro\ampere} with its supply at \SI{0}{\volt} and either I/O up to \SI{3.6}{\volt} & UQFN-10, $2\times1.5$~\si{\milli\meter} \\
\hline
\des{D}{USBESD} & ESD protection array & USBLC6-2 & IEC 61000-4-2 level 4; $\leq\SI{4}{\pico\farad}$ per line with matched line-to-line capacitance; two data lines; rail pin to VDDD; \SI{5.25}{\volt} standoff & SOT-666 or SOT23-6L \\
\hline
\des{D}{BUSTVS} & Transient suppressor & SD15Q & unidirectional; standoff $>\SI{12.6}{\volt}$; IEC 61000-4-2 level 4 stated explicitly, not a 10/1000~\si{\micro\second} surge rating alone; clamping voltage at the ESD waveform below \des{Q}{BUS}'s $V_{DS}$; capacitance unconstrained & SOD-323 or smaller \\
\hline
\des{R}{BUSG} & Film resistor & \SI{1}{\kilo\ohm} & $\pm5\,\%$ sufficient; carries the driver's full pull-down and sets \des{D}{BUSG}'s clamp current & 0402 \\
\hline
\des{R}{BUSOFF} & Film resistor & \SI{49.9}{\kilo\ohm} & $\pm5\,\%$ sufficient; gate to common source; holds the pair off with the gate undriven & 0402 \\
\hline
\des{R}{BUSBIAS} & Film resistor & \SI{10}{\kilo\ohm} & $\pm5\,\%$ sufficient; gate to the receptacle's VBUS; one series with \des{R}{BUSOFF} across a body-diode drop & 0402 \\
\hline
\des{D}{BUSG} & Zener diode & \SI{10}{\volt} & cathode on the common source, anode on the gate; clamps $V_{GS}$ against excursions beyond this rail; $\geq\SI{200}{\milli\watt}$ & SOD-323 or smaller \\
\hline
\des{C}{BUSSLEW} & Ceramic & \SI{1}{\micro\farad} nominal & $\geq\SI{25}{\volt}$; X5R or X7R; $\geq\SI{130}{\nano\farad}$ effective at \SI{12.6}{\volt}; approximately \SI{200}{\nano\farad} expected; gate to $V_{\mathrm{BUS}}$; use driver-strength setting 1 & 0402 \\
\hline
\des{C}{BUSCM} & Ceramic & \SI{1}{\micro\farad} nominal & $\geq\SI{25}{\volt}$; X5R or X7R; approximately \SI{200}{\nano\farad} expected at operating bias; gate to the receptacle's VBUS & 0402 \\
\hline
\des{R}{PDMAX1} & Film resistor & \SI{5.1}{\kilo\ohm} & 1\,\%; one series with \des{R}{PDMAX2} so the ratio tracks; the \SI{12}{\volt} window spans \SIrange{787}{1347}{\milli\volt} on the pin & 0402 \\
\hline
\des{R}{PDMAX2} & Film resistor & \SI{2.4}{\kilo\ohm} & as \des{R}{PDMAX1} & 0402 \\
\hline
\des{R}{PDI1} & Film resistor & \SI{5.1}{\kilo\ohm} & 1\,\%; one series with \des{R}{PDI2}; the $+\SI{500}{\milli\ampere}$ window spans \SIrange{787}{1347}{\milli\volt} on the pin & 0402 \\
\hline
\des{R}{PDI2} & Film resistor & \SI{2.4}{\kilo\ohm} & as \des{R}{PDI1} & 0402 \\
\hline
\des{C}{PDD1} & Ceramic, VDDD bulk & \SI{1}{\micro\farad} & $\geq\SI{10}{\volt}$; X5R or better, as the datasheet requires; \SI{0.8}{\micro\farad} effective minimum & 0603 \\
\hline
\des{C}{PDD2} & Ceramic, VDDD & \SI{100}{\nano\farad} & $\geq\SI{10}{\volt}$; X7R & 0402 \\
\hline
\des{C}{PDD3} & Ceramic, VDDD & \SI{100}{\nano\farad} & as \des{C}{PDD2} & 0402 \\
\hline
\des{C}{PDC} & Ceramic, VCCD & \SI{1}{\micro\farad} & $\geq\SI{6.3}{\volt}$; X5R or better; no other connection on the pin & 0603 \\
\hline
\des{C}{PDIN} & Ceramic, VBUS\_IN\_DIS & \SI{100}{\nano\farad} & $\geq\SI{50}{\volt}$, sized to \des{D}{BUSTVS}'s clamp and not to the rail; X5R or better; \SI{0.1}{\micro\farad} effective minimum at \SI{12.6}{\volt} bias & 0603 \\
\hline
\des{R}{PDSDA} & Film resistor & \SI{10}{\kilo\ohm} & $\pm5\,\%$ sufficient; open-drain bus pull-up to +3V3 & 0402 \\
\hline
\des{R}{PDSCL} & Film resistor & \SI{10}{\kilo\ohm} & as \des{R}{PDSDA} & 0402 \\
\hline
\des{R}{PDINT} & Film resistor & \SI{10}{\kilo\ohm} & as \des{R}{PDSDA}; the value the datasheet recommends & 0402 \\
\hline
\des{R}{PDFLT} & Film resistor & \SI{100}{\kilo\ohm} & $\pm5\,\%$ sufficient; weak pull-down selects strong-drive \SI{3.3}{\volt} mode & 0402 \\
\hline
\des{R}{USBSW} & Film resistor & \SI{10}{\kilo\ohm} & $\pm5\,\%$ sufficient; pull-up to +3V3, the switch's enable input disconnecting on a logic high; holds the switch disabled against a high-impedance MCU pin & 0402 \\
\hline
\des{R}{PDDP} & Film resistor & \SI{0}{\ohm} & joins the controller's D+ to the pair; a placement at which the controller's branch can be opened & 0402 \\
\hline
\des{R}{PDDM} & Film resistor & \SI{0}{\ohm} & as \des{R}{PDDP}, on D$-$ & 0402 \\
\hline
\des{R}{BUSBLEED} & Film resistor & \SI{100}{\kilo\ohm} & $\pm5\,\%$ sufficient; \SI{1.59}{\milli\watt} at the top of the window; populated, the controller's discharge being detach-timed and not continuous & 0402 \\
\hline
\end{longtable}
}
